%=============================================================================!
% 9.4  SPLITTING THE LIOUVILLE OPERATOR                                       !
%=============================================================================!
\section{Splitting the Liouville operator}
\label{sec:in-splitting}

Velocity Verlet was derived above from Taylor series. A second derivation,
from the Liouville operator of \cref{sec:ha-brackets}, explains why it is
built of kicks and drifts, why each of these moves is exact, and how other
methods are made, including the thermostats of Chapter~13.

\subsection*{The motion as an exponential}

\Cref{sec:ha-brackets} found that, over a short time, any quantity $A$
of the state changes along the motion as
\begin{equation*}
   A(t) = \sum_{n=0}^{\infty}\frac{t^n}{n!}\,\Liou^nA ,
\end{equation*}
with $\Liou A = \pb{A}{H}$, every term evaluated at the start. This is the
series of the exponential, $\mathrm{e}^s = \sum s^n/n!$ (\cref{sec:ne-drag}),
with $t\Liou$ in place of $s$, and it is written $\mathrm{e}^{t\Liou}A$:
the exponential of an operator is defined by this series, and applying it
to $A$ carries $A$ forward by the time $t$. For $N$ atoms with $H =
\sum_i\lvert\vb{p}_i\rvert^2/(2m_i) + U(\vb{r}^N)$, Hamilton's equations
give $\dot{\vb{r}}_i = \vb{p}_i/m_i$ and $\dot{\vb{p}}_i = \vb{F}_i$, so
\begin{equation*}
   \Liou A = \sum_i\Bigl(\frac{\vb{p}_i}{m_i}\cdot\grad_{\vb{r}_i}A
   + \vb{F}_i\cdot\grad_{\vb{p}_i}A\Bigr) = \Liou_{\mathrm r}A
   + \Liou_{\mathrm p}A ,
\end{equation*}
where $\grad_{\vb{p}_i}$ is the gradient with respect to the momentum of
atom $i$, $\Liou_{\mathrm r}$ is the first sum, the part that moves the
positions, and $\Liou_{\mathrm p}$ the second, the part that moves the
momenta. Many books write the Liouville operator $\mathrm{i}L$, with a
factor $\mathrm{i}$ taken from quantum mechanics; nothing here depends on
it.

\subsection*{Each part alone is exact}

Applied to a position, $\Liou_{\mathrm r}\vb{r}_i = \vb{p}_i/m_i$, and
applied again, $\Liou_{\mathrm r}(\vb{p}_i/m_i) = 0$, since $\vb{p}_i/m_i$
does not depend on any position. The series stops after two terms:
\begin{equation*}
   \mathrm{e}^{\dt\Liou_{\mathrm r}}\vb{r}_i = \vb{r}_i + \dt\,\frac{\vb{p}_i}
   {m_i} , \qquad \mathrm{e}^{\dt\Liou_{\mathrm r}}\vb{p}_i = \vb{p}_i ,
\end{equation*}
the second because $\Liou_{\mathrm r}\vb{p}_i = 0$. This is an exact
drift. In the same way $\Liou_{\mathrm p}\vb{p}_i = \vb{F}_i(\vb{r}^N)$,
and $\Liou_{\mathrm p}\vb{F}_i = 0$, since the forces depend only on the
positions, so
\begin{equation*}
   \mathrm{e}^{\dt\Liou_{\mathrm p}}\vb{p}_i = \vb{p}_i + \dt\,\vb{F}_i ,
   \qquad \mathrm{e}^{\dt\Liou_{\mathrm p}}\vb{r}_i = \vb{r}_i :
\end{equation*}
an exact kick. Each part of the motion alone can be followed without
error, and the error of a method built from them comes only from
combining them. Since $\Liou_{\mathrm p}A = \pb{A}{U}$ and $\Liou_{\mathrm
r}A = \pb{A}{K}$, they are the Liouville operators of $U$ alone and of $K$
alone, so by \cref{eq:ha-series} $\mathrm{e}^{\dt\Liou_{\mathrm p}}$ turns
any quantity into its value after a kick, and $\mathrm{e}^{\dt\Liou_{\mathrm
r}}$ into its value after a drift.

\subsection*{Putting the parts together}

If $\Liou_{\mathrm r}$ and $\Liou_{\mathrm p}$ behaved like numbers, then
$\mathrm{e}^{\dt(\Liou_{\mathrm r} + \Liou_{\mathrm p})}$ would equal
$\mathrm{e}^{\dt\Liou_{\mathrm r}}\,\mathrm{e}^{\dt\Liou_{\mathrm p}}$. They
do not: applying $\Liou_{\mathrm r}$ and then $\Liou_{\mathrm p}$ is not
the same as the other order, because moving the positions changes the
forces. For one coordinate, $\Liou_{\mathrm p}$ applied to $p$ gives $F$,
and $\Liou_{\mathrm r}$ applied to that gives $(p/m)\,F'(q)$, while
$\Liou_{\mathrm r}$ applied to $p$ gives 0, and $\Liou_{\mathrm p}$ applied
to that gives 0. We write $X = \Liou_{\mathrm p}$ and $Y = \Liou_{\mathrm
r}$ for short; $XY$ means $Y$ applied first and then $X$, so the order of
every product must be kept, and $1$ is the operator that leaves a quantity
unchanged. Expanding each exponential to second order in $\dt$,
\begin{align*}
   \mathrm{e}^{\dt(X + Y)} &= 1 + \dt\,(X + Y) + \tfrac12\dt^2(X + Y)^2
   + \order{\dt^3}\\
   &= 1 + \dt\,(X + Y) + \tfrac12\dt^2\bigl(X^2 + XY + YX + Y^2\bigr)
   + \order{\dt^3} ,
\end{align*}
where $(X + Y)^2 = (X + Y)(X + Y)$ is multiplied out keeping $XY$ and $YX$
apart. The symmetric product of half a kick, a drift and half a kick is
\begin{align*}
   \mathrm{e}^{\frac12\dt X}\,\mathrm{e}^{\dt Y}\,\mathrm{e}^{\frac12\dt X}
   &= \bigl(1 + \tfrac12\dt X + \tfrac18\dt^2X^2\bigr)
   \bigl(1 + \dt Y + \tfrac12\dt^2Y^2\bigr)\\
   &\qquad\times\bigl(1 + \tfrac12\dt X + \tfrac18\dt^2X^2\bigr)
   + \order{\dt^3} .
\end{align*}
Multiplying out, the terms of first order are $\frac12\dt X + \dt Y +
\frac12\dt X = \dt\,(X + Y)$, and the terms of second order are
$\frac18\dt^2X^2$ and $\frac18\dt^2X^2$ from the outer factors alone,
$\frac12\dt^2Y^2$ from the middle one, and the products of first-order
terms from two factors, $\frac12\dt X\,\dt Y$, $\dt Y\,\frac12\dt X$ and
$\frac12\dt X\,\frac12\dt X$, in the order in which the factors stand:
\begin{equation*}
   \tfrac18X^2 + \tfrac18X^2 + \tfrac12Y^2 + \tfrac12XY + \tfrac12YX
   + \tfrac14X^2 = \tfrac12\bigl(X^2 + XY + YX + Y^2\bigr) ,
\end{equation*}
times $\dt^2$. The two agree through the second order:
\begin{equation}
   \mathrm{e}^{\dt\Liou} = \mathrm{e}^{\frac12\dt\Liou_{\mathrm p}}\,
   \mathrm{e}^{\dt\Liou_{\mathrm r}}\,\mathrm{e}^{\frac12\dt\Liou_{\mathrm p}}
   + \order{\dt^3} .
   \label{eq:in-trotter}
\end{equation}
The right side, applied to the state, is a half kick, a drift and a half
kick: velocity Verlet, with the local error of order $\dt^3$ found in
\cref{sec:in-verlet}. Splitting the exponential of a sum into a symmetric
product of exponentials in this way is called a Trotter factorisation, and
Tuckerman develops this derivation in full\cite{Tuckerman2023}. In a product
of such exponentials the one written last acts first on the quantity, and
the one written first acts first on the state: in
$\mathrm{e}^{\dt\Liou_{\mathrm p}}\mathrm{e}^{\dt\Liou_{\mathrm r}}A$,
the drift turns $A(\vb{r}, \vb{p})$ into $A(\vb{r} + \dt\,\vb{p}/m,
\vb{p})$, and the kick then replaces $\vb{p}$ in this by $\vb{p} + \dt\,
\vb{F}(\vb{r})$ everywhere it appears, so the state is kicked first and
then drifts with the new momentum (\cref{ex:in-order}). That unsymmetric
product agrees with $\mathrm{e}^{\dt\Liou}$ only to first order
(\cref{ex:in-unsymmetric}), and is the symplectic Euler method of
\cref{sec:ha-symplectic}; the symmetric product reads the same in either
order.

The same idea builds other methods. Splitting the forces themselves into
fast and slow parts gives methods that kick with the slow forces less
often\cite{Tuckerman2023} (Chapter~11), and Chapter~13 adds to $\Liou_{\mathrm r}$ and $\Liou_{\mathrm p}$ a
third part, the action of a thermostat, and splits all three.

\begin{keyidea}
The motion over a step is $\mathrm{e}^{\dt\Liou}$, and the Liouville
operator splits into a part that moves the positions and a part that moves
the momenta, each of whose exponentials is exact: a drift and a kick. The
symmetric product of half a kick, a drift and half a kick agrees with
$\mathrm{e}^{\dt\Liou}$ to second order in $\dt$, and is velocity Verlet.
\end{keyidea}

\notebook{09}{the exponentials of the Liouville operator in SymPy}

\begin{selfcheck}
For one coordinate with $H = p^2/(2m)$, apply $\mathrm{e}^{\dt\Liou}$ to
$q^2$ and check that it gives $(q + \dt\,p/m)^2$.
\end{selfcheck}
